%%%%%%%%%%%%%%%%%%%%%%%%%%%%%%%%%%%%%%%%%%%%%%%%%%%%%%%%%%%%%%%%%%%
%%% Documento LaTeX 																						%%%
%%%%%%%%%%%%%%%%%%%%%%%%%%%%%%%%%%%%%%%%%%%%%%%%%%%%%%%%%%%%%%%%%%%
% Título:		Capítulo 1
% Autor:  	Xiangmin Fan
% Fecha:  	2026-03-10
%%%%%%%%%%%%%%%%%%%%%%%%%%%%%%%%%%%%%%%%%%%%%%%%%%%%%%%%%%%%%%%%%%%
% !TEX root = A0.MiTFG.tex

\chapterbegin{Introducción y visión general}
\minitoc

\begin{sinopsis}
\label{sec:intro:sinop}

El presente trabajo aborda el diseño e implementación de un procesador basado en la arquitectura RISC-V, con énfasis en el conjunto de instrucciones RV32I. El proyecto se organiza en tres fases principales. En primer lugar, se desarrolla el procesador a nivel de diseño de hardware, considerando tanto la descripción funcional como la organización interna de sus componentes mediante el uso del lenguaje de descripción de hardware Verilog. Posteriormente, se emplea un método de verificación sustentado en la generación y comparación de firmas, lo cual permite validar la corrección del comportamiento frente a la especificación del conjunto de instrucciones. Finalmente, el procesador se implementa en una plataforma FPGA, específicamente en una Digilent Basys 3, con el propósito de evaluar su desempeño en un entorno experimental y comprobar la fidelidad de la implementación.

Este trabajo contribuye a la comprensión práctica de arquitecturas abiertas como RISC-V, además de sentar una base sólida para futuras extensiones y optimizaciones en el ámbito del diseño de procesadores.

\end{sinopsis}

\section{El auge del hardware de código abierto}

El campo de la arquitectura de computadoras está experimentando una profunda transformación, liderada por el movimiento del hardware de código abierto. Durante mucho tiempo, la arquitectura del conjunto de instrucciones (ISA), que sirve como la interfaz central entre el software y el hardware, ha sido controlada de manera propietaria por un pequeño grupo de empresas, como las arquitecturas x86 y ARM. Sin embargo, la aparición de la ISA RISC-V ha roto este paradigma. Como un estándar ISA gratuito y abierto, RISC-V ha ganado rápidamente una amplia atención y adopción tanto en el mundo académico como en la industria, gracias a sus características de modularidad, extensibilidad y ausencia de regalías. No solo ha reducido la barrera de entrada para el diseño de procesadores, sino que también ha impulsado una ola de innovación sin precedentes, haciendo posible el diseño de procesadores especializados y a medida. Este proyecto elige RISC-V como base de diseño con el objetivo de explorar y poner en práctica las ventajas de ingeniería que ofrece este estándar abierto, y de contribuir con un caso de implementación concreto al ecosistema del hardware de código abierto.\cite{wiki:riscv,Free_isa}


\section{Objetivo}
\label{sec:intro:obj}

El objetivo central de este proyecto es documentar de manera integral y sistemática el proceso completo, desde la concepción hasta la implementación física, de un procesador de un solo ciclo que soporta el conjunto de instrucciones básicas de enteros RV32I. Específicamente, el proyecto busca lograr lo siguiente:

\begin{itemize}
\item Diseño e Implementación: Diseñar e implementar una CPU RISC-V de 32 bits, capaz de ejecutar correctamente todas las instrucciones RV32I. Para ello, se utilizará el lenguaje de descripción de hardware Verilog y se basará en el modelo clásico de procesador de un solo ciclo.

\item Simulación y Verificación: Establecer un flujo de verificación completo que combine la cadena de herramientas oficial de RISC-V con un entorno de simulación Verilog. Esto garantizará la correcta funcionalidad del procesador a nivel lógico.

\item Despliegue y Evaluación en FPGA: Desplegar exitosamente el diseño del procesador, una vez verificado, en una placa de desarrollo FPGA Basys 3. Se realizará una confirmación de su funcionalidad a nivel de hardware y se llevará a cabo un análisis cuantitativo del consumo de recursos y el rendimiento temporal del diseño sintetizado.

\end{itemize}


\section{Estructura del documento}

La estructura de este informe es la siguiente: la segunda parte profundiza en los fundamentos teóricos del proyecto, es decir, la arquitectura del conjunto de instrucciones RISC-V RV32I, incluyendo su modelo de programador, formato de instrucción y principios de codificación. La tercera parte detalla el diseño de la microarquitectura del procesador de un solo ciclo, cubriendo la construcción de la ruta de datos y la implementación lógica de la unidad de control. La cuarta parte muestra cómo se utiliza Verilog para traducir el diseño de la microarquitectura en código de nivel de transferencia de registros (RTL) concreto. La quinta parte presenta la estrategia de verificación del diseño, incluyendo la escritura de programas de prueba, el flujo de compilación y el análisis de los resultados de la simulación. La sexta parte describe el proceso completo de despliegue del diseño en la FPGA Basys 3, incluyendo la síntesis, la implementación y las pruebas a nivel de placa. La séptima parte analiza los resultados finales de la implementación en la FPGA, incluyendo la utilización de recursos y el rendimiento de la temporización, y presenta una visión de trabajos futuros. Finalmente, las partes octava y novena proporcionan las referencias y los apéndices, respectivamente, donde el apéndice contiene el código fuente completo del proyecto.

\section{Ámbito de aplicación}

El proyecto se enmarca dentro del campo de la Arquitectura de Computadoras y el Diseño de Sistemas Digitales. Su alcance abarca desde la conceptualización teórica hasta la validación práctica de un procesador. La implementación se centra en un diseño de un solo ciclo que ejecuta el conjunto de instrucciones RV32I de RISC-V.

El trabajo se limita al diseño de la unidad de procesamiento central (CPU), pero incluye la implementación de periféricos básicos como controladores para manejar la comunicación con dispositivos de entrada/salida (I/O). La verificación se realiza a través de la simulación en software y la validación final se lleva a cabo mediante el despliegue en un entorno FPGA, específicamente en una placa Basys 3.

Este proyecto no aborda diseños más complejos, como arquitecturas de tubería (pipelining) o procesadores multi-ciclo, sino que se enfoca en proporcionar una base sólida y comprensible de los principios fundamentales de la arquitectura de CPU. Los resultados obtenidos son aplicables a la enseñanza, la investigación académica y el desarrollo de prototipos de hardware de baja complejidad.

\chapterend
